\documentclass[12pt,a4paper]{article}
\usepackage[utf8]{inputenc}
\usepackage[spanish]{babel}
\usepackage{geometry}
\usepackage{enumitem}
\usepackage{amssymb}
\usepackage{hyperref}
\usepackage{xcolor}

\geometry{a4paper, margin=1in}

\title{Guía para la Creación de Pull Requests (PR) en GitHub \\ \large Curso: Seguridad Informática}
\author{Prof. César Rodríguez}
\date{\today}

\begin{document}

\maketitle

\section*{Introducción}
Como parte de la metodología de trabajo colaborativo en el desarrollo de nuestra Suite de Auditoría de Seguridad Informática, utilizaremos \textbf{GitHub} para integrar los aportes de cada grupo. Esta guía detalla los pasos estandarizados para enviar sus contribuciones al repositorio principal del curso mediante \textit{Pull Requests} (PR).

\section{Paso 1: Realizar un \textit{Fork} del Repositorio}
Dado que los estudiantes no tienen permisos de escritura directa en el repositorio \texttt{main} del curso, el primer paso es crear una copia personal en su propia cuenta de GitHub.
\begin{enumerate}
    \item Ingrese a la página del repositorio principal del curso en GitHub.
    \item Haga clic en el botón \textbf{'Fork'} ubicado en la esquina superior derecha.
    \item Seleccione su cuenta personal como destino para copiar el repositorio.
\end{enumerate}

\section{Paso 2: Clonar su Repositorio Localmente}
Abra su terminal y clone el repositorio que ahora reside en su cuenta personal. Reemplace \texttt{SU\_USUARIO} con su nombre de usuario de GitHub:
\begin{verbatim}
git clone https://github.com/SU_USUARIO/CYBERSEGURIDAD.git
cd CYBERSEGURIDAD
\end{verbatim}

\section{Paso 3: Crear una Rama de Trabajo (\textit{Branch})}
Nunca trabaje directamente sobre la rama \texttt{main}. Para mantener el orden, cree una rama específica para la tarea asignada a su grupo (ej. DNS, OSINT, Discovery, Scanning):
\begin{verbatim}
git checkout -b feature/grupoX-estudianteY
\end{verbatim}

\section{Paso 4: Realizar los Cambios y Hacer \textit{Push}}
Trabaje \textbf{únicamente} en el archivo asignado a su grupo dentro de la carpeta \texttt{modulos/} (ej. \texttt{modulos/scanning.py}). Recuerde la \textbf{Regla de Oro}: ¡Está estrictamente prohibido modificar \texttt{auditoria.py} o los archivos de otros grupos!

Una vez finalizada y probada su función localmente:
\begin{verbatim}
git add modulos/su_archivo_de_grupo.py
git commit -m 'Implementada funcion XYZ segun los requerimientos'
git push origin feature/grupoX-estudianteY
\end{verbatim}

\section{Paso 5: Crear el \textit{Pull Request}}
\begin{enumerate}
    \item Vaya a la página de su repositorio (\textit{fork}) en GitHub desde su navegador web.
    \item Verá un banner verde que dice \textbf{'Compare \& pull request'}. Haga clic en él.
    \item Asegúrese de que la rama base (\textit{base repository/branch}) sea la del curso y la rama de comparación (\textit{head repository/branch}) sea la de su cuenta.
    \item En la descripción del PR, \textbf{debe} incluir una explicación clara de lo que desarrolló, mencionando su Grupo y su Número de Estudiante.
    \item Haga clic en \textbf{'Create pull request'}.
\end{enumerate}

\section*{Lista de Verificación antes de Enviar}
Para que su Pull Request sea aprobado e integrado, asegúrese de cumplir con los siguientes puntos:
\begin{itemize}
    \item[\checkmark] \textbf{Contrato de Datos:} Su función retorna un diccionario válido según \texttt{docs/schema\_resultados.json}.
    \item[\checkmark] \textbf{Calidad de Código:} Incluyó \textit{Type Hinting} y \textit{Docstrings} (Estilo Google).
    \item[\checkmark] \textbf{Resiliencia:} Su bloque de código está protegido por \texttt{try-except} para capturar errores de red sin detener la auditoría.
    \item[\checkmark] \textbf{Limpieza:} Eliminó la instrucción \texttt{raise NotImplementedError} de la plantilla base.
\end{itemize}

\vspace{0.5cm}
\noindent \textit{Nota: Si el profesor realiza comentarios o solicita cambios en su PR, no necesita crear uno nuevo. Simplemente haga los cambios en su computadora, haga \texttt{commit} y \texttt{push} nuevamente a su rama. El PR se actualizará de forma automática.}

\end{document}